\section{习题}

\subsection{本章速览}
\begin{itemize}
    \item 级数 $\sum_{n=1}^{\infty}a_n$ 收敛的必要条件是 $a_n\to0$，但这不是充分条件；判断敛散性归根到底是看部分和 $S_n$ 是否有极限。
    \item 正项级数单调递增的部分和只要有上界就收敛。常用工具：比较判别法（及其极限形式）、比值判别法、根值判别法、积分判别法；记住 $p$ 级数 $\sum n^{-p}$ 当且仅当 $p>1$ 时收敛。
    \item 比值/根值判别法：若 $\rho<1$ 则收敛，若 $\rho>1$ 则发散，若 $\rho=1$ 则失效，需要另想办法。
    \item 交错级数用 Leibniz 判别法：$u_n$ 单调不增且 $u_n\to0$，则 $\sum(-1)^{n-1}u_n$ 收敛，且余项满足 $|R_n|\le u_{n+1}$。绝对收敛必收敛；绝对收敛与条件收敛要分开讨论。
    \item 幂级数 $\sum a_n(x-x_0)^n$ 的收敛半径由 $R=1/\lim|a_{n+1}/a_n|$ 或 $R=1/\lim\sqrt[n]{|a_n|}$ 给出，端点必须单独讨论。在收敛区间内可逐项求导、逐项积分，收敛半径不变。
    \item 傅里叶系数为 $a_n=\frac1l\int_{-l}^{l}f(x)\cos\frac{n\pi x}{l}\,dx$，$b_n=\frac1l\int_{-l}^{l}f(x)\sin\frac{n\pi x}{l}\,dx$。若 $f$ 满足狄利克雷条件，则其傅里叶级数在连续点处收敛于 $f(x)$，在第一类间断点处收敛于左右极限的平均值。
\end{itemize}

\subsection{基础过关题}

\textbf{1.} 用比值判别法或根值判别法判定下列正项级数的敛散性：
(1) $\displaystyle\sum_{n=1}^{\infty}\frac{n!}{3^n}$；
(2) $\displaystyle\sum_{n=1}^{\infty}\frac{n^2}{2^n}$；
(3) $\displaystyle\sum_{n=1}^{\infty}\left(\frac{n}{2n+1}\right)^n$。

\ifshowsolutions%
\boxed{\textbf{解}}
\textbf{思路：} 通项含 $n!$ 时，相邻两项相除能大量约分，优先用比值判别法；通项本身是 $n$ 次幂时，开 $n$ 次方最方便，优先用根值判别法。

(1) 记 $a_n=\dfrac{n!}{3^n}$，则
\[
\frac{a_{n+1}}{a_n}
=\frac{(n+1)!/3^{n+1}}{n!/3^n}
=\frac{n+1}{3}\longrightarrow+\infty>1.
\]
比值极限大于 $1$，故级数发散。

(2) 记 $a_n=\dfrac{n^2}{2^n}$，则
\[
\frac{a_{n+1}}{a_n}
=\frac{(n+1)^2/2^{n+1}}{n^2/2^n}
=\frac12\left(1+\frac1n\right)^2\longrightarrow\frac12<1.
\]
故级数收敛。

(3) 记 $a_n=\left(\dfrac{n}{2n+1}\right)^n$，则
\[
\sqrt[n]{a_n}=\frac{n}{2n+1}\longrightarrow\frac12<1.
\]
故级数收敛。

\boxed{\parbox{0.9\linewidth}{\textbf{答案：}(1) 发散；(2) 收敛；(3) 收敛。}}
\else \answerarea[8] \fi%

\textbf{2.} 判定下列正项级数的敛散性：
(1) $\displaystyle\sum_{n=1}^{\infty}\sin\frac{1}{n^2}$；
(2) $\displaystyle\sum_{n=2}^{\infty}\frac{1}{n\ln^2 n}$；
(3) $\displaystyle\sum_{n=1}^{\infty}\ln\left(1+\frac1n\right)$。

\ifshowsolutions%
\boxed{\textbf{解}}
\textbf{思路：} 先看通项趋于零的速度，再与已知的 $p$ 级数或调和级数比较；第 (2) 题是积分判别法的典型形式。

(1) 当 $n\to\infty$ 时 $\dfrac{1}{n^2}\to0$，故
\[
\lim_{n\to\infty}\frac{\sin(1/n^2)}{1/n^2}=1.
\]
而 $\sum\dfrac{1}{n^2}$ 是 $p=2>1$ 的 $p$ 级数，收敛，所以原级数收敛。

(2) 令 $f(x)=\dfrac{1}{x\ln^2 x}$，它在 $[2,+\infty)$ 上非负、连续且单调递减，并且 $f(n)=\dfrac{1}{n\ln^2 n}$。由积分判别法，
\[
\int_2^{+\infty}\frac{dx}{x\ln^2 x}
=\lim_{b\to+\infty}\left[-\frac{1}{\ln x}\right]_2^b
=\frac{1}{\ln2}<+\infty.
\]
积分收敛，故原级数收敛。

(3) 由等价无穷小 $\ln(1+t)\sim t\ (t\to0)$，取 $t=\dfrac1n$，得
\[
\lim_{n\to\infty}\frac{\ln(1+1/n)}{1/n}=1.
\]
而调和级数 $\sum\dfrac1n$ 发散，故原级数发散。

\boxed{\parbox{0.9\linewidth}{\textbf{答案：}(1) 收敛；(2) 收敛；(3) 发散。}}
\else \answerarea[8] \fi%

\textbf{3.} 判断级数 $\displaystyle\sum_{n=1}^{\infty}(-1)^{n-1}\frac{n}{n^2+1}$ 是绝对收敛、条件收敛还是发散。

\ifshowsolutions%
\boxed{\textbf{解}}
\textbf{思路：} 先去掉符号判断绝对值级数；若绝对值级数发散，再用 Leibniz 判别法研究原交错级数。

令 $u_n=\dfrac{n}{n^2+1}$。先看绝对值级数：
\[
\frac{n}{n^2+1}\ge\frac{1}{2n}\qquad(n\ge1),
\]
因为 $2n^2\ge n^2+1$ 对 $n\ge1$ 成立。而 $\sum\dfrac{1}{2n}$ 发散，由比较判别法，$\sum\dfrac{n}{n^2+1}$ 发散。

再对原级数用 Leibniz 判别法。令 $f(x)=\dfrac{x}{x^2+1}$，则
\[
f'(x)=\frac{1-x^2}{(x^2+1)^2}\le0\qquad(x\ge1),
\]
所以 $\{u_n\}$ 从 $n=1$ 起单调不增。又
\[
\lim_{n\to\infty}\frac{n}{n^2+1}=0.
\]
Leibniz 判别法的两个条件都满足，故原交错级数收敛。

由于绝对值级数发散而原级数收敛，原级数为条件收敛。

\boxed{\parbox{0.9\linewidth}{\textbf{答案：}条件收敛。}}
\else \answerarea[8] \fi%

\textbf{4.} 判断下列命题是否正确，并说明理由：
(1) 若 $\sum a_n$ 收敛，则 $\sum a_n^2$ 收敛；
(2) 若 $\sum a_n$ 收敛，则 $\sum(-1)^na_n$ 收敛；
(3) 若 $\sum a_n$ 绝对收敛，则 $\sum a_n^2$ 收敛；
(4) 若 $\sum a_n$ 条件收敛，则 $\sum a_n^2$ 收敛。

\ifshowsolutions%
\boxed{\textbf{解}}
\textbf{思路：} 收敛只保证 $a_n\to0$，并不能保证 $|a_n|$ 可求和；构造“条件收敛但平方后发散”的例子即可推翻 (1)(2)(4)。

取
\[
a_n=\frac{(-1)^n}{\sqrt n}.
\]
由 Leibniz 判别法，$\sum a_n=\sum\dfrac{(-1)^n}{\sqrt n}$ 收敛；但
\[
\sum|a_n|=\sum\frac{1}{\sqrt n}
\]
是 $p=\dfrac12\le1$ 的 $p$ 级数，发散，所以 $\sum a_n$ 只是条件收敛。

(1) 对上述 $a_n$，有 $a_n^2=\dfrac1n$，故 $\sum a_n^2$ 是调和级数，发散。命题错误。

(2) 对上述 $a_n$，
\[
\sum(-1)^na_n=\sum\frac{1}{\sqrt n},
\]
发散。命题错误。

(3) 若 $\sum a_n$ 绝对收敛，则 $a_n\to0$，故存在 $N$，当 $n\ge N$ 时 $|a_n|\le1$。于是
\[
a_n^2=|a_n|^2\le|a_n|\qquad(n\ge N).
\]
由比较判别法，$\sum a_n^2$ 收敛。命题正确。

(4) 对上述 $a_n$，$\sum a_n$ 条件收敛，但 $\sum a_n^2=\sum\dfrac1n$ 发散。命题错误。

\boxed{\parbox{0.9\linewidth}{\textbf{答案：}只有 (3) 正确；(1)(2)(4) 均错误。}}
\else \answerarea[8] \fi%

\textbf{5.} 求幂级数 $\displaystyle\sum_{n=1}^{\infty}\frac{(x-2)^n}{n\,3^n}$ 的收敛域。

\ifshowsolutions%
\boxed{\textbf{解}}
\textbf{思路：} 先由系数求出收敛半径，得到开区间；再分别把两个端点代回原级数判断。

记 $a_n=\dfrac{1}{n3^n}$，则
\[
\lim_{n\to\infty}\left|\frac{a_{n+1}}{a_n}\right|
=\lim_{n\to\infty}\frac{n3^n}{(n+1)3^{n+1}}
=\frac13.
\]
故收敛半径
\[
R=\frac{1}{1/3}=3,
\]
收敛区间为 $|x-2|<3$，即 $-1<x<5$。

当 $x=5$ 时，级数为 $\displaystyle\sum_{n=1}^{\infty}\frac{1}{n}$，是调和级数，发散。

当 $x=-1$ 时，级数为 $\displaystyle\sum_{n=1}^{\infty}\frac{(-1)^n}{n}$，由 Leibniz 判别法收敛。

因此收敛域为 $[-1,5)$。

\boxed{\parbox{0.9\linewidth}{\textbf{答案：}$[-1,5)$。}}
\else \answerarea[8] \fi%

\subsection{综合提高题}

\textbf{6.} 求幂级数 $\displaystyle S(x)=\sum_{n=1}^{\infty}n^2x^n$ 的收敛域与和函数，并求数项级数 $\displaystyle\sum_{n=1}^{\infty}\frac{n^2}{2^n}$ 的和。

\ifshowsolutions%
\boxed{\textbf{解}}
\textbf{思路：} 先由几何级数得到 $\sum nx^{n-1}$，再乘 $x$、求导，得到 $n^2x^n$ 的和函数。

记 $a_n=n^2$，则
\[
\lim_{n\to\infty}\left|\frac{a_{n+1}}{a_n}\right|
=\lim_{n\to\infty}\left(1+\frac1n\right)^2=1,
\]
故收敛半径 $R=1$。当 $|x|=1$ 时，通项 $n^2x^n$ 不趋于零，级数发散，所以收敛域为 $(-1,1)$。

在 $|x|<1$ 内，由几何级数
\[
\sum_{n=0}^{\infty}x^n=\frac{1}{1-x}
\]
逐项求导得
\[
\sum_{n=1}^{\infty}nx^{n-1}=\frac{1}{(1-x)^2}.
\]
两边乘以 $x$：
\[
\sum_{n=1}^{\infty}nx^n=\frac{x}{(1-x)^2}.
\]
再逐项求导：
\[
\sum_{n=1}^{\infty}n^2x^{n-1}
=\frac{d}{dx}\left[\frac{x}{(1-x)^2}\right]
=\frac{(1-x)^2+2x(1-x)}{(1-x)^4}
=\frac{1+x}{(1-x)^3}.
\]
两边乘以 $x$，得
\[
S(x)=\sum_{n=1}^{\infty}n^2x^n
=\frac{x(1+x)}{(1-x)^3},\qquad |x|<1.
\]
取 $x=\dfrac12$：
\[
\sum_{n=1}^{\infty}\frac{n^2}{2^n}
=S\left(\frac12\right)
=\frac{\frac12\cdot\frac32}{\left(\frac12\right)^3}
=\frac{3/4}{1/8}=6.
\]

\boxed{\parbox{0.9\linewidth}{\textbf{答案：}收敛域为 $(-1,1)$，$S(x)=\dfrac{x(1+x)}{(1-x)^3}$；$\displaystyle\sum_{n=1}^{\infty}\frac{n^2}{2^n}=6$。}}
\else \answerarea[8] \fi%

\textbf{7.} 求幂级数 $\displaystyle\sum_{n=1}^{\infty}\frac{(x-1)^{2n}}{n\,4^n}$ 的收敛域与和函数。

\ifshowsolutions%
\boxed{\textbf{解}}
\textbf{思路：} 把 $(x-1)^2$ 整体看作一个新变量，转化为标准的 $\sum t^n/n$。

令
\[
t=\frac{(x-1)^2}{4}\ge0.
\]
原级数化为 $\displaystyle\sum_{n=1}^{\infty}\frac{t^n}{n}$。该级数在 $0\le t<1$ 时收敛，在 $t\ge1$ 时发散（$t=1$ 时为调和级数，$t>1$ 时通项不趋于零）。

由 $t<1$ 得
\[
\frac{(x-1)^2}{4}<1
\quad\Longleftrightarrow\quad
|x-1|<2,
\]
即 $-1<x<3$。当 $x=-1$ 或 $x=3$ 时 $t=1$，级数发散。故收敛域为 $(-1,3)$。

在 $|x-1|<2$ 内，由
\[
\sum_{n=1}^{\infty}\frac{t^n}{n}=-\ln(1-t)\qquad(0\le t<1),
\]
得
\[
S(x)=-\ln\left(1-\frac{(x-1)^2}{4}\right)
=\ln\frac{4}{4-(x-1)^2}
=\ln\frac{4}{3+2x-x^2}.
\]

\boxed{\parbox{0.9\linewidth}{\textbf{答案：}收敛域为 $(-1,3)$，$S(x)=\ln\dfrac{4}{3+2x-x^2}$。}}
\else \answerarea[8] \fi%

\textbf{8.} 求下列数项级数的和：
(1) $\displaystyle\sum_{n=1}^{\infty}\frac{n^2}{n!}$；
(2) $\displaystyle\sum_{n=1}^{\infty}(-1)^n\frac{n^2}{n!}$。

\ifshowsolutions%
\boxed{\textbf{解}}
\textbf{思路：} 利用 $n^2=n(n-1)+n$ 拆分，使分子与阶乘约分，再用 $e^x$ 的展开式。

因为 $\sum\dfrac{n^2}{n!}$ 与 $\sum(-1)^n\dfrac{n^2}{n!}$ 都绝对收敛（通项与 $\dfrac{1}{(n-2)!}$ 同阶），可以拆项。

(1) 由 $n^2=n(n-1)+n$，
\[
\sum_{n=1}^{\infty}\frac{n^2}{n!}
=\sum_{n=2}^{\infty}\frac{1}{(n-2)!}
+\sum_{n=1}^{\infty}\frac{1}{(n-1)!}.
\]
第一个级数令 $k=n-2$，第二个令 $k=n-1$，得
\[
\sum_{k=0}^{\infty}\frac{1}{k!}+\sum_{k=0}^{\infty}\frac{1}{k!}
=e+e=2e.
\]

(2) 同理，
\[
\sum_{n=1}^{\infty}(-1)^n\frac{n^2}{n!}
=\sum_{n=2}^{\infty}\frac{(-1)^n}{(n-2)!}
+\sum_{n=1}^{\infty}\frac{(-1)^n}{(n-1)!}.
\]
第一项令 $k=n-2$：
\[
\sum_{k=0}^{\infty}\frac{(-1)^{k+2}}{k!}
=\sum_{k=0}^{\infty}\frac{(-1)^k}{k!}=e^{-1}.
\]
第二项令 $k=n-1$：
\[
\sum_{k=0}^{\infty}\frac{(-1)^{k+1}}{k!}
=-\sum_{k=0}^{\infty}\frac{(-1)^k}{k!}=-e^{-1}.
\]
两项相加得 $e^{-1}-e^{-1}=0$。

\boxed{\parbox{0.9\linewidth}{\textbf{答案：}(1) $2e$；(2) $0$。}}
\else \answerarea[8] \fi%

\textbf{9.} 将 $f(x)=x\ (-\pi<x<\pi)$ 展开为以 $2\pi$ 为周期的傅里叶级数，并利用展开式求
\[
\sum_{k=1}^{\infty}\frac{(-1)^{k-1}}{2k-1}.
\]

\ifshowsolutions%
\boxed{\textbf{解}}
\textbf{思路：} $f(x)=x$ 是奇函数，所以余弦项系数全为零，只需算正弦项系数；再在连续点 $x=\dfrac{\pi}{2}$ 处代入。

因为 $f(x)=x$ 是奇函数，所以
\[
a_0=0,\qquad a_n=0\quad(n\ge1).
\]
而
\[
b_n=\frac1\pi\int_{-\pi}^{\pi}x\sin nx\,dx
=\frac2\pi\int_0^{\pi}x\sin nx\,dx.
\]
计算
\[
\int x\sin nx\,dx=-\frac{x\cos nx}{n}+\frac{\sin nx}{n^2},
\]
故
\[
\int_0^{\pi}x\sin nx\,dx
=\left[-\frac{x\cos nx}{n}+\frac{\sin nx}{n^2}\right]_0^{\pi}
=-\frac{\pi(-1)^n}{n}.
\]
于是
\[
b_n=\frac2\pi\left(-\frac{\pi(-1)^n}{n}\right)
=\frac{2(-1)^{n+1}}{n}.
\]
因此
\[
x=2\sum_{n=1}^{\infty}\frac{(-1)^{n+1}}{n}\sin nx,
\qquad -\pi<x<\pi.
\]

取 $x=\dfrac{\pi}{2}$，此时 $f$ 连续，级数收敛于 $\dfrac{\pi}{2}$：
\[
\frac{\pi}{2}
=2\sum_{n=1}^{\infty}\frac{(-1)^{n+1}}{n}\sin\frac{n\pi}{2}.
\]
当 $n$ 为偶数时 $\sin\dfrac{n\pi}{2}=0$；当 $n=2k-1$ 时，
\[
\sin\frac{(2k-1)\pi}{2}=(-1)^{k-1},\qquad
(-1)^{n+1}=(-1)^{2k}=1.
\]
所以
\[
\frac{\pi}{2}
=2\sum_{k=1}^{\infty}\frac{(-1)^{k-1}}{2k-1},
\]
即
\[
\sum_{k=1}^{\infty}\frac{(-1)^{k-1}}{2k-1}=\frac{\pi}{4}.
\]

\boxed{\parbox{0.9\linewidth}{\textbf{答案：}$x=2\displaystyle\sum_{n=1}^{\infty}\frac{(-1)^{n+1}}{n}\sin nx\ (-\pi<x<\pi)$；$\displaystyle\sum_{k=1}^{\infty}\frac{(-1)^{k-1}}{2k-1}=\frac{\pi}{4}$。}}
\else \answerarea[10] \fi%

\subsection{真题风格与综合训练}

\textbf{10.}（真题风格）设 $a>0$，讨论级数
\[
\sum_{n=1}^{\infty}\frac{a^n n!}{n^n}
\]
的敛散性。

\ifshowsolutions%
\boxed{\textbf{解}}
\textbf{思路：} 通项含 $n!$ 和 $n^n$，先用比值判别法；当比值极限等于 $1$ 时判别法失效，此时用 Stirling 公式估计通项。

记 $u_n=\dfrac{a^n n!}{n^n}>0$，则
\[
\frac{u_{n+1}}{u_n}
=\frac{a^{n+1}(n+1)!}{(n+1)^{n+1}}\cdot\frac{n^n}{a^n n!}
=a\left(\frac{n}{n+1}\right)^n.
\]
而
\[
\left(\frac{n}{n+1}\right)^n
=\frac{1}{(1+1/n)^n}\longrightarrow\frac1e,
\]
所以
\[
\lim_{n\to\infty}\frac{u_{n+1}}{u_n}=\frac ae.
\]

当 $0<a<e$ 时，$\dfrac ae<1$，由比值判别法，级数收敛。

当 $a>e$ 时，$\dfrac ae>1$，比值判别法表明级数发散。

当 $a=e$ 时，比值极限为 $1$，比值判别法失效。但此时
\[
\frac{u_{n+1}}{u_n}
=e\left(\frac{n}{n+1}\right)^n
=\frac{e}{(1+1/n)^n}.
\]
而数列 $\left(1+\frac1n\right)^n$ 单调递增且趋于 $e$，故对每个 $n$ 都有
\[
\left(1+\frac1n\right)^n<e,
\]
从而 $\dfrac{u_{n+1}}{u_n}>1$，即 $\{u_n\}$ 严格递增。又 $u_1=e>0$，所以 $u_n\ge e$，通项不趋于零，级数发散。

\boxed{\parbox{0.9\linewidth}{\textbf{答案：}当 $0<a<e$ 时收敛；当 $a\ge e$ 时发散。}}
\else \answerarea[10] \fi%

\textbf{11.}（真题风格）将 $f(x)=|\sin x|$ 在 $[-\pi,\pi]$ 上展开为傅里叶级数，并利用展开式求
\[
\sum_{k=1}^{\infty}\frac{1}{4k^2-1}.
\]

\ifshowsolutions%
\boxed{\textbf{解}}
\textbf{思路：}$|\sin x|$ 是偶函数，因此 $b_n=0$，只算余弦系数。积分时用积化和差，并注意 $n=1$ 要单独处理。

因为 $f(x)=|\sin x|$ 是偶函数，所以
\[
b_n=0\quad(n\ge1).
\]
常数项系数
\[
a_0=\frac1\pi\int_{-\pi}^{\pi}|\sin x|\,dx
=\frac2\pi\int_0^{\pi}\sin x\,dx
=\frac2\pi\cdot2=\frac4\pi.
\]

当 $n=1$ 时，
\[
a_1=\frac2\pi\int_0^{\pi}\sin x\cos x\,dx
=\frac2\pi\int_0^{\pi}\frac12\sin2x\,dx=0.
\]
当 $n\ge2$ 时，由积化和差
\[
\sin x\cos nx=\frac12\left[\sin(n+1)x-\sin(n-1)x\right],
\]
以及 $\displaystyle\int_0^{\pi}\sin mx\,dx=\frac{1-(-1)^m}{m}\ (m\ge1)$，得
\[
a_n=\frac2\pi\cdot\frac12
\left[
\frac{1-(-1)^{n+1}}{n+1}
-\frac{1-(-1)^{n-1}}{n-1}
\right].
\]
注意到 $(-1)^{n+1}=(-1)^{n-1}$。当 $n$ 为偶数时，这两个数都等于 $-1$，故
\[
a_n=\frac1\pi\left(\frac{2}{n+1}-\frac{2}{n-1}\right)
=-\frac{4}{\pi(n^2-1)}.
\]
当 $n\ge3$ 为奇数时，这两个数都等于 $1$，故 $a_n=0$。

$|\sin x|$ 在 $[-\pi,\pi]$ 上连续且满足狄利克雷条件，因此
\[
|\sin x|
=\frac{2}{\pi}-\frac{4}{\pi}\sum_{k=1}^{\infty}\frac{\cos 2kx}{4k^2-1},
\qquad -\pi\le x\le\pi.
\]

令 $x=0$，左边为 $0$，右边各余弦项均为 $1$，于是
\[
0=\frac2\pi-\frac4\pi\sum_{k=1}^{\infty}\frac{1}{4k^2-1},
\]
解得
\[
\sum_{k=1}^{\infty}\frac{1}{4k^2-1}=\frac12.
\]

\boxed{\parbox{0.9\linewidth}{\textbf{答案：}$|\sin x|=\dfrac2\pi-\dfrac4\pi\displaystyle\sum_{k=1}^{\infty}\frac{\cos 2kx}{4k^2-1}$；$\displaystyle\sum_{k=1}^{\infty}\frac{1}{4k^2-1}=\frac12$。}}
\else \answerarea[10] \fi%

\textbf{12.}（真题风格）求幂级数
\[
S(x)=\sum_{n=1}^{\infty}\frac{x^n}{n(n+2)}
\]
的收敛域与和函数，并求 $\displaystyle\sum_{n=1}^{\infty}\frac{1}{n(n+2)2^n}$。

\ifshowsolutions%
\boxed{\textbf{解}}
\textbf{思路：} 先求收敛半径并检查端点；求和时把 $\dfrac{1}{n(n+2)}$ 拆成部分分式，化为两个对数级数。

记 $a_n=\dfrac{1}{n(n+2)}$，则
\[
\lim_{n\to\infty}\left|\frac{a_{n+1}}{a_n}\right|
=\lim_{n\to\infty}\frac{n(n+2)}{(n+1)(n+3)}=1,
\]
故收敛半径 $R=1$。当 $x=\pm1$ 时，
\[
\left|\frac{x^n}{n(n+2)}\right|=\frac{1}{n(n+2)}\le\frac{1}{n^2},
\]
而 $\sum\dfrac1{n^2}$ 收敛，所以两个端点都绝对收敛。收敛域为 $[-1,1]$。

当 $0<|x|<1$ 时，由部分分式
\[
\frac{1}{n(n+2)}=\frac12\left(\frac1n-\frac1{n+2}\right),
\]
得
\[
S(x)=\frac12\left[
\sum_{n=1}^{\infty}\frac{x^n}{n}
-\sum_{n=1}^{\infty}\frac{x^n}{n+2}
\right].
\]
其中
\[
\sum_{n=1}^{\infty}\frac{x^n}{n}=-\ln(1-x).
\]
又
\[
\sum_{n=1}^{\infty}\frac{x^n}{n+2}
=\frac1{x^2}\sum_{n=1}^{\infty}\frac{x^{n+2}}{n+2}
=\frac1{x^2}\left[-\ln(1-x)-x-\frac{x^2}{2}\right].
\]
代入得
\[
S(x)=\frac12\left[
-\ln(1-x)+\frac{\ln(1-x)}{x^2}+\frac1x+\frac12
\right]
=\frac12\left[
\frac12+\frac1x+\left(\frac1{x^2}-1\right)\ln(1-x)
\right].
\]
另外 $S(0)=0$。在端点处由和函数的连续性（或直接计算）得
\[
S(1)=\sum_{n=1}^{\infty}\frac{1}{n(n+2)}
=\frac12\left(1+\frac12\right)=\frac34.
\]
又
\[
S(-1)=\frac12\left[\sum_{n=1}^{\infty}\frac{(-1)^n}{n}-\sum_{n=1}^{\infty}\frac{(-1)^n}{n+2}\right]
=\frac12\left[-\ln2-\left(-\ln2+\frac12\right)\right]
=-\frac14,
\]
其中第二个级数令 $k=n+2$，利用 $\sum_{k=1}^{\infty}\frac{(-1)^k}{k}=-\ln2$ 及前两项 $k=1,2$ 的和 $-1+\frac12=-\frac12$ 得到。

取 $x=\dfrac12$：
\[
\sum_{n=1}^{\infty}\frac{1}{n(n+2)2^n}
=S\left(\frac12\right)
=\frac12\left[\frac12+2+3\ln\frac12\right]
=\frac54-\frac32\ln2.
\]

\boxed{\parbox{0.9\linewidth}{\textbf{答案：}收敛域为 $[-1,1]$；当 $0<|x|<1$ 时 $S(x)=\frac12\left[\frac12+\frac1x+\left(\frac1{x^2}-1\right)\ln(1-x)\right]$，且 $S(0)=0$，$S(1)=\frac34$，$S(-1)=-\frac14$；$\displaystyle\sum_{n=1}^{\infty}\frac{1}{n(n+2)2^n}=\frac54-\frac32\ln2$。}}
\else \answerarea[10] \fi%
